% Nouveau EGA V (cote 25): archive cote 25, site page 45 (PDF page 46).
% typescript. Draft compared: none.
% AI-transcribed by Claude Opus 5.5 (claude-opus-5-5, adaptive thinking, effort high), Message Batches API, batch ega25b (msgbatch_01A7n5YuXubUPNK1zsepi8RG), 2026-11-13; original AI transcription, subsequently checked against both supplied scan parts; see PAGES_CHECKED.tsv and CHECK_DECISIONS.tsv.
\scanpage{25/45}
% NOTE: The small left-margin blot has uncertain function. The detached acute-like pen stroke over lesquels is retained without assigning grammatical force. The circled insertion places sing on a lower baseline than sup.

\mspage{\hand{41}}

hg$_1$ :  $g^{-1}(U) \hand{\longrightarrow}$ PxT  n'est autre que l'inclusion canonique, puisqu'il
en est ainsi de sa restriction à $g^{-1}(V) \hand{\rightarrow}$ V.  \var{\xout{Soit alors \unsure{U'|U} = h$^{-1}$(}}{}{typewriter strike-out}{ega25_25_045-v0001}
\var{\xout{Comme h prend}}{}{typewriter strike-out}{ega25_25_045-v0002} Il en résulte que h \var{\xout{prend ses valeurs dans}}{}{typewriter strike-out}{ega25_25_045-v0003} se factorise
à travers l'adhérence schématique \var{\xout{de}}{}{typewriter strike-out}{ega25_25_045-v0004} $\overset{\hand{\rule{1.1em}{0.35pt}}}{Y^1}$ de Y$^1$ dans  PxT , donc que
l'image inverse de Y$^1$ (qui est un ouvert de ladite adhérence) par h
est un ouvert de U, soit U' . En vertu de  h g$_1$ = inclusion, on
voit alors aussitôt que  g$_1$ induit un isomorphisme  $g^{-1}(U) \hand{\simeq}$ U', et
\var{\xout{comme}}{}{typewriter strike-out}{ega25_25_045-v0005} cela montre bien que  $g^{-1}(U) \hand{\longrightarrow}$ U est une immersion ouverte.
Lorsque X est propre sur k i.e. f ermé dans P, alors  g est propre donc
surjectif, donc  g$_1$ : $g^{-1}(U) \hand{\longrightarrow}$ U est surjectif, donc un isomorphisme.
D'ailleurs, si  W est un ouvert de T tel que  $g^{-1}(W) \hand{\longrightarrow}$ W soit un isomor%
phisme, il s'ensuit que W est lisse puisque Y$^1$ est lisse, donc W$\hand{\subset}$U.
Cela achève la démonstration de 8.\var{\over{\unsure{8}}{9}}{9}{writing over}{ega25_25_045-v0006}. 

\var{\xout{Remarque}}{}{typewriter strike-out}{ega25_25_045-v0007}  L\var{\over{'}{es}}{es}{writing over}{ega25_25_045-v0008} assertion\hand{s} finale\hand{s} de 8.\var{\over{\unsure{7}}{8}}{8}{writing over}{ega25_25_045-v0009}.\hand{,} \var{\inserted{$\underline{Y}^{\text{sup}}\text{sing}=\emptyset$ ou}}{$\underline{Y}^{\text{sup}}\text{sing}=\emptyset$ ou}{insertion}{ega25_25_045-v0010} dim $\underline{Y}^{\text{sup sing}}$ \var{\over{=}{\unsure{$\leq$}}}{\unsure{$\leq$}}{writing over}{ega25_25_045-v0011} r-2\hand{,} et
dim T = r-1 , sont triviales;  la première résult\var{\over{ant}{e}}{e}{writing over}{ega25_25_045-v0012} d\var{\over{u}{\unsure{e}}}{\unsure{e}}{writing over}{ega25_25_045-v0013} \var{\deleted{$\underline{Y}^{\text{sup sing}} \neq \underline{Y}^{\text{sing}}$}}{}{deletion}{ega25_25_045-v0014}
\var{\inserted{\var{\deleted{$\underline{Y}^{\text{sup sing}}$}}{}{deletion}{ega25_25_045-v0016}}}{\var{\deleted{$\underline{Y}^{\text{sup sing}}$}}{}{deletion}{ega25_25_045-v0016}}{insertion}{ega25_25_045-v0015} \var{\inserted{\var{\deleted{contenu}}{}{deletion}{ega25_25_045-v0018} fait que $\underline{Y}^{\text{sing}}$, de dim r}}{\var{\deleted{contenu}}{}{deletion}{ega25_25_045-v0018} fait que $\underline{Y}^{\text{sing}}$, de dim r}{insertion}{ega25_25_045-v0017} \var{\deleted{donc rare dans $\underline{Y}^{\text{sing}}$ puisque celui}}{}{deletion}{ega25_25_045-v0019} est irréductible, et du fait que
\var{\deleted{dim}}{}{deletion}{ega25_25_045-v0020} $\underline{Y}^{\var{\over{\text{sing}}{\text{sup sing}}}{\text{sup sing}}{writing over}{ega25_25_045-v0021}}$ \var{\deleted{= r-1}}{}{deletion}{ega25_25_045-v0022} \var{\inserted{est défini par l'annulation d'une section D d'un Module inversible}}{est défini par l'annulation d'une section D d'un Module inversible}{insertion}{ega25_25_045-v0023} ; la deuxième du fait que, L étant fini sur K, on a
deg tr L/k = deg\var{\xout{x}}{}{typewriter strike-out}{ega25_25_045-v0024} tr K/k, i.e.  \var{\xout{dim $\underline{Y}^{\text{sing}}$ xxxxxx}}{}{typewriter strike-out}{ega25_25_045-v0025} dim T = dim $\underline{Y}^{\text{sing}}$ = r-1.

\uline{Remarque} 8.\var{\over{\unsure{9}}{10}}{10}{writing over}{ega25_25_045-v0026}.  Comme on a si\var{\over{n}{g}}{g}{writing over}{ega25_25_045-v0027}\var{\over{g}{n}}{n}{writing over}{ega25_25_045-v0028}alé dans  8.\var{\over{\unsure{8}}{9}}{9}{writing over}{ega25_25_045-v0029}. , avec les notations de
ce corollaire, on a  $g^{-1}(U) \hand{\subset} \underline{Y}^{\text{sing}} - \underline{Y}^{\text{sup sing}}$ ; mais on notera que
même si X est fermé dans P,  cette inclusion n'est pas nécessairement
\var{\margin{\hand{\unsure{$\cdot$}}}}{\placedmargin{\hand{\unsure{$\cdot$}}}}{marginal note placed}{ega25_25_045-v0030} une égalité, en d'autres termes (notant que $g^{-1}(U)$ n'est autre que
l'ensemble des points où g est \uline{étale}, alors que $\underline{Y}^{\text{sing}}$-$\underline{Y}^{\text{sup sing}}$ est
lèensemble des points où g est \uline{non ramifié}) il peut y avoir des points
y de $\underline{Y}^{\text{sing}}$ en \hand{´}lesquels  g est \var{\xout{étale}}{}{typewriter strike-out}{ega25_25_045-v0031} non ramifié, mais non étale (ce qui
\var{\inserted{\var{\deleted{\illegible{1}}}{}{deletion}{ega25_25_045-v0033}}}{\var{\deleted{\illegible{1}}}{}{deletion}{ega25_25_045-v0033}}{insertion}{ega25_25_045-v0032} implique d'ailleurs que g(y) est un point non \var{\xout{normal}}{}{typewriter strike-out}{ega25_25_045-v0034} géom normal et même
non géom unibranche de  T)\hand{.} \var{\xout{Soit \illegible{1}}}{}{typewriter strike-out}{ega25_25_045-v0035} En termes géométriques, cela corres%
pond au phénomène suivant: on peut avoir un hyperplan tangent non oscul
% NOTE: the last line is cut at the right edge of the scan
\var{\margin{\hand{45}}}{\placedmargin{\hand{45}}}{marginal note placed}{ega25_25_045-v0036}
